\begin{center}
\begin{tikzpicture}[x=.9cm,y=.9cm,font=\small]
  \fill[paper] (0,0) rectangle (10.5,5.7);
  \draw[blue!55,line width=1pt] (.3,4.7) .. controls (3.0,4.4) and (6.0,5.0) .. (10.0,4.4);
  \node[text=blue!70!black,above] at (5.2,4.65) {黄河};
  \fill[orange!22] (1.0,3.0) -- (4.0,3.0) -- (4.4,5.1) -- (1.5,5.2) -- cycle;
  \node[align=center] at (2.6,4.05) {燕\\北方边缘};
  \fill[cyan!22] (6.9,1.2) -- (10.0,1.3) -- (10.1,4.0) -- (7.3,4.1) -- cycle;
  \node[align=center] at (8.6,2.7) {齐\\临淄方向};
  \fill[timeline] (5.4,3.4) circle (.14);
  \node[above,align=center] at (5.4,3.5) {济西\\战场方向};
  \draw[-{Stealth[length=2mm]},ultra thick,color=orange!70!black] (3.6,3.9) -- (5.25,3.45);
  \draw[-{Stealth[length=2mm]},ultra thick,color=orange!70!black] (5.55,3.4) -- (7.0,2.9);
  \node[text=orange!70!black,below] at (5.2,2.95) {燕与五国联军东进};
\end{tikzpicture}

\smallskip
\small\textit{图  燕伐齐空间示意：突出燕的北方位置、齐的东方腹地、济西方向与联军东进；不表示精确战线或五国军队的实际部署。}
\end{center}
